% 第 28 章　热力学第一定律：能量总账
\chapter{热力学第一定律：能量总账}

\step{反常识开场}

\begin{mainline}
找一个自行车打气筒，给轮胎快速打上二三十下，然后蹲下来摸一摸金属筒身——尤其是靠近出气口的那一段。它是热的，有时热得明显。

再换一个场景。夏天用按压式的喷雾罐——杀虫剂、降温喷雾、喷漆都行——按住喷头持续喷十几秒，然后摸一摸铁罐。它是凉的，凉到罐壁上常常能凝出一层水珠。

把这两个画面并排摆在一起，怪事就出现了。打气筒没有靠近过任何火源，房间里没有任何比它更热的东西，它凭什么变热？喷雾罐没有进过冰箱，周围都是暖洋洋的空气，它凭什么变凉？我们从小熟悉的经验是“热从热的东西流向冷的东西”，可这两件事里，热既不像是从哪儿流来的，冷也不像是被谁送来的。热和冷仿佛是被凭空“做”出来的。

类似的怪事其实到处都是：搓几秒钟手，手心就热了；汽车急刹一次，刹车片能烫到冒烟；钻木可以取火。反过来，气球放气时气球嘴冰凉，高压锅喷气时喷出的蒸汽迅速变凉。热能被凭空制造，冷也能被凭空制造——如果热是一种“东西”，这东西难道不守恒吗？这一章，我们要给能量建一本总账，把这笔看似混乱的账目记清楚。
\end{mainline}

\step{先猜一猜}

\begin{mainline}
往下读之前，先押注。请你对下面三个问题各给一个直觉答案，写下来：

\begin{itemize}
  \item 打气筒发热，是不是活塞和筒壁之间的摩擦“磨”出来的？
  \item 把气球的气一口气放掉，气球嘴附近会变热还是变凉？
  \item 一杯热水放在桌上慢慢变凉，散掉的能量去了哪里？这个过程能不能自发地倒过来——一杯凉咖啡从空气里吸热，自己重新沸腾？
\end{itemize}

先剧透一句：第一题绝大多数人都会答错，而且错得很有价值——因为错误的答案里藏着一个被高估了两百年的旧模型。第二题你可能凭经验答对，但未必说得出理由。第三题的两小问，答案分别是“说得出”和“说不出”，而“说不出”的那部分，恰恰是本章最后要留给下一章的钩子。
\end{mainline}

\step{动手实验}

\begin{experiment}[打气筒与一根橡皮筋]
你需要：一个自行车打气筒（金属筒身最好），一根宽橡皮筋（扎饭盒、捆蔬菜的那种），你的嘴唇——它是你身上最灵敏的温度计。

第一步，打气筒。给轮胎连续快速打气二三十下，停下来立刻摸筒身的不同高度。你会发现靠近出气口的底部最热，越往上越不明显。然后做一个关键的对照实验：把气管从气嘴上拔下来，让活塞空打——空气自由进出、几乎不被压缩——用同样的频率、同样的次数再打二三十下，再摸筒身。这一次明显不怎么热。活塞和筒壁的摩擦次数一模一样，发热却少了大半。记住这个反差，它是对“摩擦生热解释打气筒”的第一记重拳。

第二步，橡皮筋。把橡皮筋松松地贴在上唇上，快速把它拉长到原来的两三倍，保持住——你会清楚地感到它在发热。保持拉伸十几秒，让它把热散掉，然后快速松手让它缩回去，立刻再贴嘴唇——它变凉了。

第三步，停下来想。两个实验有一个共同的模式：变热的那一步，都有人对东西“使劲”——压缩空气、拉长橡皮筋；变凉的那一步，都是东西自己“缩回去”——膨胀、回缩。使劲与发热、回缩与变凉之间的这条对应关系，就是本章公式要翻译的直觉。

安全提示：橡皮筋不要对着自己或别人的眼睛拉；打气筒连续使用会很热，摸之前先用手背轻触。
\end{experiment}

\step{旧模型遇到麻烦}

\begin{mainline}
在牛顿力学里，我们已经有一本能量账：动能加势能，叫作机械能，在只有重力、弹力这类力做功时守恒。拿这本旧账来记本章的现象，会遇到三个过不去的麻烦。

麻烦一：刹车的账销不掉。一辆汽车急刹停下，几十万千焦的动能不见了。它既没有变成高度（车还在原地），也没有变成弹簧的压缩量。旧账本上，这笔钱只能记成“消失了”。搓手、钻木、打气筒也一样：机械能明明投了进去，却没有以任何动能或势能的形式存下来——东西只是变热了。如果能量账本允许这样不明不白地销账，那它根本就不配叫守恒定律。

麻烦二：旧的热模型解释不了“热被造出来”。十八世纪的主流模型把热想象成一种没有质量的流体，叫作“热质”：热的物体热质多，冷的物体热质少，接触时热质从多的地方流向少的地方。这个模型能漂亮地解释热水变凉、冷水变温之类的现象。但 1798 年，伦福德伯爵在慕尼黑的兵工厂里监督镗钻大炮炮筒时发现：用钝钻头钻金属，产生的金属碎屑很少，而炮筒周围水槽里的水却能一直热下去——只要钻头继续转，水就能一直沸腾。如果热是从金属里“挤”出来的热质，一块金属里储存的热质怎么会取之不尽？摩擦生热看起来根本没有上限。

麻烦三：打气筒现场根本没有“更热的物体”。你的手是三十七摄氏度左右，筒身却可以热到超过体温。如果热只能从高温流向低温，筒身的热是从谁那里流来的？热质说在这里彻底破产：热不是在搬家，而是在被生产。

几乎同一时期，德国医生迈尔从另一条路摸到了同一个结论。他在远航船上当随船医生时注意到，船员在热带抽出的静脉血比在寒带更红——他推断，热带气温高，身体为维持体温需要“燃烧”的食物更少，血液里剩下的氧就更多。食物、体热、劳力之间似乎存在某种固定的兑换关系。一个医生从血液颜色想到能量守恒，听上去离奇，却说明时机已经成熟：只缺一个人把它测准。

这个人就是焦耳。他在十九世纪四十年代做了一系列实验，最出名的装置是：让重物匀速下落，通过绳子和滑轮带动浸在水里的桨叶搅动，然后用温度计读水的温升。结果干净利落：让一克水升高一摄氏度所需的热，总是精确地等价于约 4.2 焦耳的机械功——不管这份功来自搅动、压缩还是摩擦。机械能和热之间存在一个固定不变的“汇率”。这意味着热不是一种物质，而是能量的一种存在和转移方式。刹车的账有救了：动能没有消失，它换了一个科目存了起来。这个新科目，就是下一节的主角。
\end{mainline}

\step{发明新概念}

\begin{mainline}
要像焦耳那一代人那样重建账本，先得发明三个记账用的概念：\textbf{系统}、\textbf{环境}、\textbf{边界}。

系统，是我们决定跟踪的那一小部分世界：一瓶气体、一缸水、一台发动机。环境，是系统之外、可能与系统交换能量的一切。边界，是两者之间的分界——它可以是真实的器壁，也可以只是你想象中的一个面；它可以固定不动，也可以跟着活塞移动。这三个词听起来平淡，却是热力学的全部语法：账记不记错，几乎全看边界画得清不清楚。每当你对一道热学题感到糊涂，先问自己：我的系统是谁？边界画在哪？

接下来是账本上的三个科目。

第一笔：\textbf{内能}，记作 \(U\)。它是系统“此刻拥有”的能量总和：分子无规运动的动能、分子之间相互作用的势能、锁在化学键里的能量，都计入其中（微观图像见第 26、27 章）。内能最重要的身份是\textbf{状态量}：它只由系统此刻的状态决定——多少物质、什么温度、多大体积——而与系统“怎么走到这个状态”完全无关。一瓶气体无论你先压缩后加热，还是先加热后压缩，只要最后到达同一个状态，内能就是同一个数。绕一圈回到出发状态，内能必定恢复原值。

第二笔和第三笔：\textbf{热量} \(Q\) 与\textbf{功} \(W\)。它们和内能的地位截然不同——它们不是系统“含有”的东西，而是能量\textbf{穿过边界}的两种方式，叫作过程量。功，是通过宏观的、有组织的动作推过边界的能量：活塞整体移动一段距离、转轴拧过一个角度、电流驱动一台电机，微观上看是亿万分子被“整齐划一”地推了一把。热，是仅靠温度差、通过边界上无数次无规碰撞“漏”过去的能量，没有任何宏观位移。所以“这瓶气体含有多少热”是一句病句——热只在跨越边界的那一刻才有意义，正如你不能问一个银行账户“余额里有多少是转账”。

这里值得把银行比喻说完整。它像的地方：内能像账户余额，只看现在、不问来历；热和功像“转账”与“现金”两种存取方式，钱的总数变化等于两者之和。它不像的地方有两处。其一，银行里转账记录可以事后查询、可以撤销，而能量一旦进入内能，就永远无法分辨哪一焦耳来自热、哪一焦耳来自功——这正是“过程量”的深意。其二，也是本章末尾的关键：银行流水可以倒着走，而能量账对应的真实过程有方向，第一定律的账本上偏偏没有记录这条规则。记住这个缺口。

顺带一提，这套新语言恰好接上了全书反复追问的三个问题。系统此刻是什么状态？由温度、压强、体积等少数几个量刻画，内能是状态的函数。什么规律决定状态如何变化？第一定律是答案的一半——能量怎么过账；另一半“哪些变化允许发生”留给后面的章节。我们怎样通过测量知道它？内能不能直接读表，但功可以量（力乘距离）、热可以量（温度计加比热），账目两边各自可测，平衡与否可以被实验当场检验——这正是焦耳实验做过的事。
\end{mainline}

\begin{figure}[htbp]
\centering
\begin{tikzpicture}[scale=1]
  \draw[rounded corners=6pt] (-4.6,-2.2) rectangle (4.6,2.6);
  \node at (0,2.95) {环境};
  \draw[thick,rounded corners=14pt] (-1.8,-1.1) rectangle (1.8,1.1);
  \node at (0,0) {系统（内能 \(U\)）};
  \node at (2.6,-1.5) {边界};
  \draw[->] (2.2,-1.35) -- (1.55,-0.85);
  \draw[-{Stealth},very thick] (-4.0,0.4) -- (-2.1,0.4) node[midway,above] {热量 \(Q\)（流入）};
  \draw[-{Stealth},very thick] (2.1,-0.5) -- (4.0,-0.5) node[midway,below] {功 \(W\)（对外）};
\end{tikzpicture}
\caption{系统、环境与边界。内能是系统此刻的“存款”，热量与功只在穿过边界时有意义。}
\end{figure}

\step{一句话模型}

\begin{oneline}
能量总账：系统内能的变化，等于穿过边界流入的热量，减去系统对外做的功；热和功只是两种过账方式，内能才是余额——账永远平衡，但余额永远不告诉你钱是怎么进来的。
\end{oneline}

\step{数学翻译器}

\begin{mainline}
把一句话模型写成符号，就是热力学第一定律：
\[
  \Delta U \;=\; Q \;-\; W .
\]
按老规矩，回答四个问题。

\textbf{它描述什么？} 一个边界清楚的系统经历任何过程之后，内能的变化量，与两股穿过边界的能量流之间的关系。左边 \(\Delta U\) 是“存款变动”，右边 \(Q\) 与 \(W\) 是仅有的两种存取方式。

\textbf{为什么这样写？} 因为能量不生不灭。系统内能的增加，只能来自两笔钱：以热的方式流进来的，减去以功的方式付出去的。减号纯粹来自符号约定：本书约定 \(Q>0\) 表示系统吸热，\(W>0\) 表示系统\textbf{对外}做功。有些教科书采用相反的约定（\(W>0\) 表示外界对系统做功），公式就写成 \(\Delta U=Q+W\)——物理内容完全一样，只是记账方向不同。读任何一本热学书，先确认它的约定，这是省钱的好习惯。

\textbf{何时成立？} 只要系统和边界定义清楚、所有能量形式都登记入账，从打气筒到柴油发动机到恒星，它都成立。它是物理学中适用范围最广的定律之一。

\textbf{何时失效或需要扩展？} 有三种情况要修补：系统与外界交换物质（拧开盖的瓶子）；涉及核反应、静质量本身参与能量转换；以及尺度小到量子涨落显著时。细节见“模型的边界”。

按本书惯例，公式一出场立刻用特殊情况检查一遍：

\begin{itemize}
  \item \(Q=0\)（绝热过程）：\(\Delta U=-W\)。气体膨胀对外做功，只能掏自己的存款——内能减少，温度下降。这就是喷雾罐变凉的账目。
  \item \(W=0\)（刚性容器，体积不变）：\(\Delta U=Q\)。密封容器在火上烤，吸收的热全部存为内能。
  \item \(Q=W=0\)（孤立系统）：\(\Delta U=0\)。总能量不变——原来“能量守恒”只是第一定律在孤立系统上的特例，第一定律比它更一般。
  \item 循环一周回到原状态：\(\Delta U=0\)，于是 \(Q=W\)。系统净吸收的热全部变成净输出的功。这是后面热机的全部账目基础。
  \item 方向反转：压缩气体时 \(W<0\)，于是 \(-W>0\)，内能增加。打气筒发热的账目就在这里。
\end{itemize}

接下来回答一个自然的问题：功具体怎么算？考虑活塞截面积为 \(S\)，气体压强为 \(p\)，把活塞缓慢推出一小段距离 \(\Delta x\)。气体对活塞的推力是 \(pS\)，做的功约为 \(pS\,\Delta x\)。注意 \(S\,\Delta x\) 正是体积的增加量 \(\Delta V\)，所以每一小步的功就是
\[
  W_{\text{小步}} \approx p\,\Delta V .
\]
如果压强在膨胀过程中不断变化（通常正是如此），就把整个过程切成许多小步，每步用当时的压强乘以体积增量，再全部加起来。这个“压强对体积的累加”有一个漂亮的几何形象：在以体积为横轴、压强为纵轴的图上画出过程曲线，\textbf{曲线下的面积就是功}。这是 \(p\)-\(V\) 图成为热力学第一图表的原因——它不只是一张状态地图，还是一台功的计数器。
\end{mainline}

\begin{figure}[htbp]
\centering
\begin{tikzpicture}[scale=1.15]
  \draw[-{Stealth}] (0,0) -- (6.4,0) node[below left] {体积 \(V\)};
  \draw[-{Stealth}] (0,0) -- (0,3.6) node[above] {压强 \(p\)};
  \draw[thick,domain=1.1:5.4,samples=60] plot (\x,{3.3/\x});
  \fill[pattern=north east lines] (1.6,0) -- plot[domain=1.6:4.6,samples=40] (\x,{3.3/\x}) -- (4.6,0) -- cycle;
  \draw[dashed] (1.6,0) -- (1.6,{3.3/1.6});
  \draw[dashed] (4.6,0) -- (4.6,{3.3/4.6});
  \node[below] at (1.6,0) {\(V_1\)};
  \node[below] at (4.6,0) {\(V_2\)};
  \node at (3.05,0.55) {面积 \(=W\)};
  \node[right] at (4.15,1.75) {膨胀过程曲线};
\end{tikzpicture}
\caption{\(p\)-\(V\) 图：过程曲线下的阴影面积等于气体对外做的功。膨胀时面积为正，压缩时体积减小、功为负。}
\end{figure}

\begin{mainline}
有了这台计数器，可以清点四种最常用的“招牌过程”，仍以理想气体为例：

\begin{itemize}
  \item \textbf{等温过程}：让气缸泡在恒温水浴里、缓慢移动，温度始终不变。膨胀时，吸收的热量恰好全部变成对外做的功，内能原地不动（理想气体的内能只由温度决定）。
  \item \textbf{等容过程}：把容器焊死，体积不变，\(W=0\)。加热时吸的热一分不少全部存进内能。
  \item \textbf{等压过程}：活塞上放着固定重物，压强不变。加热时吸的热兵分两路：一部分存为内能，一部分顶起活塞对外做功。
  \item \textbf{绝热过程}：把系统裹进厚棉被，或者干脆让过程快到热量来不及进出——打气筒的一个冲程就是后者。做功与内能此消彼长，一分做功一分内能。
\end{itemize}

同一份气体，走不同的路径，账目的分法完全不同；但无论你选哪条路径，第一定律的总账永远平衡。这正是状态量与过程量的分工：\(\Delta U\) 只看起点终点，\(Q\) 与 \(W\) 各看各的路径，而三者永远被 \(\Delta U=Q-W\) 拴在一起。

现在做两个带真实数字的估算，把“热功等价”从口号变成感觉。

估算一：焦耳的瀑布。水从约 110 米高的瀑布顶端落到底部，假设重力势能全部变成水的内能。每千克水释放的能量是 \(9.8\times110\approx1.1\times10^{3}\) 焦耳，而每千克水升高一摄氏度需要约 4200 焦耳，所以温升只有 \(1100/4200\approx0.26\) 摄氏度。据说焦耳在蜜月旅行时还带着温度计，想在瀑布的上游和下游之间测出这点温差——瀑布上下确实应该有温差，但不到零点三度的信号完全被水花掺混和蒸发淹没，他没有成功。思路却完全正确：机械功和热之间有一个确定的兑换率，约 4.2 焦耳换一卡。这个数叫热功当量，它就是当年“热质说”的验尸报告，也是能量总账的出生证明。

估算二：你一天的食物是多大一笔能量？成年人一天大约吃进 2000 千卡，约合 8.4 兆焦。如果把这笔能量全部用来把 70 千克的你竖直抬高，高度是 \(8.4\times10^{6}/(70\times9.8)\approx1.2\times10^{4}\) 米——超过珠穆朗玛峰。你的身体每天“烧掉”的能量，相当于把自己抬上珠峰还有富余。食物是极其浓缩的能量账户，而你的身体是一台从不平白赊账的热力学机器。

估算三：烧开一壶水，有多少能量花在“推开空气”上？一克水在一百摄氏度汽化要吸热约 2260 焦（这叫汽化潜热）；与此同时，一克水变成蒸汽后体积从约 1 立方厘米膨胀到约 1700 立方厘米，要在大气压下“腾出地方”，做的功约为 \(p\,\Delta V\approx10^{5}\times1.7\times10^{-3}\approx170\) 焦。也就是说，汽化吸的热里约九成三存进了蒸汽的内能（用来拆散水分子之间的相互吸引），只有约百分之七花在顶开大气上。电水壶烧干时房间并没有被“推大”，这份功悄无声息地混进了空气的内能里——账照样平。
\end{mainline}

\begin{mathbridge}
把“面积就是功”写成微积分语言：功是压强对体积的积分
\[
  W=\int_{V_1}^{V_2} p\,\mathrm{d}V .
\]
对理想气体，等温过程有 \(p=nRT/V\)，代入积分得
\[
  W=nRT\ln\frac{V_2}{V_1}.
\]
检查：\(V_2>V_1\) 时对数为正，膨胀对外做正功；压缩时对数为负，外界向系统存钱——符号规则自动正确。

绝热过程没有热流，第一定律给出另一条曲线：可以推出
\[
  pV^{\gamma}=\text{常数},\qquad TV^{\gamma-1}=\text{常数},
\]
其中 \(\gamma\) 叫绝热指数，空气约为 \(1.4\)。因为 \(\gamma>1\)，绝热线在 \(p\)-\(V\) 图上比等温线更陡：同样压缩到一半体积，绝热过程的压强升得更高——因为压缩功全部留在气体里变成了升温。

用这条曲线估算打气筒：一个冲程把空气压缩到约七分之一体积（对应打进约七个大气压的轮胎），若完全绝热，末态温度
\[
  T_2=300\times 7^{\,0.4}\approx 650\ \mathrm{K},
\]
接近四百摄氏度。这是把全部压缩功都留给气体本身的理论上限；实际过程中气体边被压缩边通过筒壁散热，达不到这么高，但这个上限足以解释筒身为什么烫手——筒壁正是这笔能量的“二传手”。柴油机把同一个原理用到极致：压缩比高达十六到二十二，空气被绝热压缩到燃料燃点以上，柴油喷入即自燃，根本不需要火花塞。压燃，是绝热过程最赚钱的工程应用。
\end{mathbridge}

\begin{challenge}
工程与化学中更常用的账本不是 \(U\)，而是焓 \(H=U+pV\)。对等压过程（敞口反应釜、大气压下的加热），由第一定律可直接推出 \(Q_p=\Delta H\)：等压下吸收的热量等于焓的增加，其中 \(p\,\Delta V\) 那一部分被自动记成“推开大气所付的账”。化学家测的反应热大多是 \(\Delta H\) 而不是 \(\Delta U\)，原因就在这里。进一步，对持续有物质流进流出的稳流设备（涡轮机、压缩机、喷管），第一定律的实用形式要以焓为主角并计入动能项：
\[
  \dot Q-\dot W_{\text{轴}}=\dot m\left[\left(h_2-h_1\right)+\tfrac12\left(v_2^2-v_1^2\right)\right].
\]
形式上比 \(\Delta U=Q-W\) 复杂，灵魂还是同一句话：边界画好之后，流进来的减去流出去的，等于系统里存下的。另外提醒一句本章用不到的边界情形：当过程涉及静质量与能量的转换时，内能科目里必须把 \(mc^2\) 登记入账；登记之后，总账依然平衡。
\end{challenge}

\step{模型的边界}

\begin{mainline}
第一定律的账本建立在两个前提上：系统和边界可以画清楚；所有能量形式都被登记入账。前提松动的地方，就是模型需要修补的地方。

修补一：开放系统。拧开盖的汽水瓶、持续进气的发动机气缸，物质本身在穿越边界，随物质一起搬运的内能和“把物质推进推出所需的功”都必须入账。这正是上一节挑战层里焓与稳流公式的用武之地——不是第一定律失效，而是记账科目要升级。

修补二：核过程。普通物理化学过程中原子核不动，静质量是一笔从不动的“死存款”；一旦涉及裂变、聚变或粒子湮灭，静质量本身会参与兑换，必须把 \(E=mc^2\) 计入内能。计入之后第一定律纹丝不动——科目增加了，平衡的原则没有倒。

修补三：微观尺度。对单个分子或几十个分子的小系统，\(Q\)、\(W\)、\(U\) 都在不停涨落，“哪一份算热、哪一份算功”的划分需要更仔细的定义。这不影响宏观总账，但提醒我们：热与功的区分本身是宏观概念，是统计意义上才干净的分类。

但第一定律最深的那条边界不在这些技术细节上，而在它管不到的一件事上：\textbf{方向}。回到本节开头的凉咖啡。热水放凉，热量流出、内能下降，账是平的。现在想象逆过程：一杯凉咖啡从周围空气里吸热，自己重新沸腾，空气相应地变冷。这笔账平不平？完全平——流入的 \(Q\) 等于增加的 \(\Delta U\)，第一定律毫无意见。可这样的过程从来、从来没有被任何人看到过。打碎的杯子不会自己拼回去，揉皱的纸不会自己展平，这些逆过程个个都不违反能量总账。可见，自然界除了“账必须平”之外，还执行着另一条独立的规则，决定哪些账允许发生、哪些永远不许。这条规则叫热力学第二定律，它是第 29、30 章的主角；第一定律只是记账员，不是立法者。

边界的最后一站，是一个会让直觉当场出错的反例：\textbf{热泵}。冬天用空调制热，说明书上写着“消耗 1 度电，可向室内供热相当于 3 度电”。效率百分之三百？能量守恒被打破了？都不是。热泵根本不是把电“变成”热，而是用电做功，把室外的热“搬”进室内：室内得到的热量 \(=\) 从室外取走的热量 \(+\) 输入的电功，账严丝合缝。真正“一度电变一度热”的是电暖器，效率百分之一百——而电暖器才是那个把所有电功都变成热的普通选手。热泵超出的部分全是搬运所得，不是创造所得。区分“转化”与“搬运”，是读懂一切热机、冰箱、空调的钥匙。

顺带纠一个常见的语言误用：“能量守恒，所以能源用不完”。总量守恒不等于可用性守恒。汽油烧掉后能量一点没少，只是变成了散布在空气里、再也无法轻易收拢的形式。账面上分文不差，能用的部分却一去不回——这正是“方向问题”在日常语言里的投影。
\end{mainline}

\step{回头解释开场现象}

\begin{mainline}
现在收网，用总账解释开场的一热一冷。

打气筒为什么热？压缩冲程中，活塞对被封闭的气体做功；冲程很快，热量来不及穿过筒壁散掉，气体近似经历绝热压缩。由 \(\Delta U=-W\)，压缩时 \(W<0\)，内能增加，气体温度升高；随后热从高温气体传给金属筒壁，再传到你的手上。所以筒身最热处在底部——那里是被压缩气体停留和放热的地方。对照实验也闭环了：空打时摩擦照常存在，却几乎不热，因为压缩功这一大头消失了，摩擦只是小头。“先猜一猜”第一题的错误答案，至此批改完毕。

开场提到的另一件事也能算出账来：汽车刹车。一辆 1300 千克的家用车以时速 120 千米行驶，动能约 \(72\) 万焦。一次急刹，这笔能量几乎全部变成刹车盘和刹车片的内能：假设其中一半热量由四片共约 2 千克的钢质刹车片吸收，它们的温升可达三四百摄氏度——难怪山路长下坡时刹车会冒烟、甚至暂时失灵（老司机说的“刹车热衰减”）。动能没有消失，它只是换了一个你摸得着但收不回来的科目。

喷雾罐为什么凉？罐内高压气体喷出时迅速膨胀，推开外界大气做功；喷嘴附近过程快、换热少，近似绝热。由 \(\Delta U=-W\)，膨胀时 \(W>0\)，内能减少，温度下降；罐身被罐内剩余的冷气体冷却，空气中的水汽便在罐壁上凝结成珠。真实气体里分子间的相互作用还会带来额外的冷却效应，结论的方向不变。

同一个原理在自然界有一台巨大的演示装置：焚风。潮湿气流翻越山脉，迎风坡被迫上升，空气绝热膨胀变冷（每升高一百米大约降温半度到一度），水汽凝结成雨；翻过山顶后，变干的空气沿背风坡下沉，被越来越高的大气压绝热压缩，每下降一百米升温约一度。结果背风坡又干又热——阿尔卑斯山北麓、安第斯山东侧的居民对此再熟悉不过。同一本能量账，山这边付了“凝结放热”，山那边收回“压缩升温”，第一定律全程平衡。

工程上，热机、冰箱、热泵是三兄弟，共用一张能流图。热机从高温热源吸热 \(Q_{\text{高}}\)，一部分变成对外输出的功 \(W\)，剩下的 \(Q_{\text{低}}\) 排给低温热源；循环一周 \(\Delta U=0\)，所以 \(W=Q_{\text{高}}-Q_{\text{低}}\)。蒸汽机、汽车发动机、燃气轮机、核电站的汽轮机，全是这个结构。冰箱与热泵把箭头倒过来：外界输入功 \(W\)，把热量 \(Q_{\text{低}}\) 从低温处搬到高温处，高温处收到 \(Q_{\text{高}}=Q_{\text{低}}+W\)。空调夏天制冷是冰箱模式，冬天制热是热泵模式，一台机器两个名字，账是同一本。
\end{mainline}

\begin{figure}[htbp]
\centering
\begin{tikzpicture}[scale=1]
  % 热机
  \draw[fill=red!8,rounded corners=3pt] (-4.7,1.7) rectangle (-2.3,2.4);
  \node at (-3.5,2.05) {高温热源};
  \draw[fill=blue!8,rounded corners=3pt] (-4.7,-2.4) rectangle (-2.3,-1.7);
  \node at (-3.5,-2.05) {低温热源};
  \draw[thick] (-3.5,0) circle (0.55);
  \node at (-3.5,0) {机};
  \draw[-{Stealth},very thick] (-3.5,1.7) -- (-3.5,0.6) node[midway,left] {\(Q_{\text{高}}\)};
  \draw[-{Stealth},very thick] (-3.5,-0.6) -- (-3.5,-1.7) node[midway,left] {\(Q_{\text{低}}\)};
  \draw[-{Stealth},very thick] (-2.95,0) -- (-1.7,0) node[right] {功 \(W\)};
  % 冰箱/热泵
  \draw[fill=blue!8,rounded corners=3pt] (2.3,1.7) rectangle (4.7,2.4);
  \node at (3.5,2.05) {高温处（室内）};
  \draw[fill=blue!8,rounded corners=3pt] (2.3,-2.4) rectangle (4.7,-1.7);
  \node at (3.5,-2.05) {低温处（室外）};
  \draw[thick] (3.5,0) circle (0.55);
  \node at (3.5,0) {泵};
  \draw[-{Stealth},very thick] (3.5,-1.7) -- (3.5,-0.6) node[midway,right] {\(Q_{\text{低}}\)};
  \draw[-{Stealth},very thick] (3.5,0.6) -- (3.5,1.7) node[midway,right] {\(Q_{\text{高}}\)};
  \draw[-{Stealth},very thick] (1.7,0) -- (2.95,0) node[midway,above] {功 \(W\)};
\end{tikzpicture}
\caption{左：热机从高温吸热，一部分变功，其余排向低温。右：冰箱与热泵输入功，把热从低温搬到高温。箭头方向相反，账目同样平衡。}
\end{figure}

\begin{mainline}
最后看一眼这张能流图上那个刺眼的细节：热机为什么不能把 \(Q_{\text{高}}\) 全部变成 \(W\)？为什么一定要排掉一份 \(Q_{\text{低}}\)？第一定律对此保持沉默——把热全部变成功、一点不排，账目照样平衡。可两百年来的工程师发现，无论设计多么精巧，这份“排放”永远省不掉。记账员又一次指向了立法者：什么决定了热机必须浪费、过程必须有方向？答案藏在“打碎的杯子为什么不会自己复原”里，那是第 29 章的故事；而它凝练成的热力学第二定律，在第 30 章等着我们。至于“能量为什么守恒”这个更深处的问题，要等到讨论对称性的章节才有真正的回答：能量守恒对应的，是“物理规律不随时间改变”这条对称性——账本之所以平衡，是因为时间的规则本身从不改账。
\end{mainline}

\step{脑洞实验与挑战题}

\begin{thought}[膨胀的宇宙是一台绝热机吗]
把第一定律推到它能到达的最大尺度：整个宇宙。在大尺度上，宇宙像一团持续膨胀的气体，而且它没有“外界”可以换热——按定义，宇宙之外没有环境。这是一个终极的绝热系统。第一定律于是给出预言：膨胀必须由内能买单，宇宙应当越胀越冷。

这正是真实发生的事。早期宇宙炽热致密；随着膨胀，它一路降温。宇宙微波背景辐射——大爆炸留下的“余温”——在宇宙年龄约三十八万年时对应约三千开尔文的等离子体，经过一百三十八亿年的膨胀，今天只剩下约 \(2.7\) 开尔文，恰好落在射电望远镜测到的数值上。需要坦白的是，辐射的冷却规律与气体分子不完全相同：光子的数目和能量要另算一本账，“膨胀对谁做功”这个问题最终要交给广义相对论。但“没有外界可换热，膨胀就降温”这条第一定律的骨架，一直撑到了宇宙学尺度。反过来想更有味道：如果宇宙的账可以不平，我们连这台最大机器的过去都无从推算；能把一百三十八亿年的账一路算到今天，本身就是能量总账最宏大的一次平衡审计。

它像打气筒吗？像的部分：绝热、膨胀、降温。不像的部分：没有活塞、没有筒壁，“对外做功”的对象不存在，功这个科目本身在宇宙尺度需要重新定义——这正是把旧账本用到新尺度时，必须缴纳的新学费。
\end{thought}

\begin{puzzles}
  \item 夏天太热，有人灵机一动：把冰箱门敞开，不就能给厨房降温了吗？可行吗？（思路提示：把“冰箱加厨房”划成一个系统。冰箱内部确实变凉，但它的排热管就在厨房里，排出的热量等于从箱内吸走的热量\textbf{加上}压缩机消耗的全部电能。总账：房间反而净得热量，越开越热。真正的空调之所以有效，是因为排热管挂到了窗外——区别只在边界画在哪里。）
  \item 一个绝热的刚性容器被隔板分成两半，一半充着理想气体，一半是真空。抽掉隔板，气体自由膨胀充满整个容器。气体的温度会变吗？（思路提示：绝热所以 \(Q=0\)；没有活塞可推、真空不需要付功，所以 \(W=0\)；于是 \(\Delta U=0\)。理想气体的内能只由温度决定，故温度不变。直觉喊“膨胀就会变冷”，是把“推动活塞的绝热膨胀”错套了过来——自由膨胀里没有任何对象收到这份功。追问：真实气体严格不变吗？分子间有相互作用势能，温度会略微改变——这正是历史上测量分子间作用力的经典实验思路。）
  \item 两份完全相同的气体从同一状态出发，一份等温压缩、一份绝热压缩，压到相同的更小体积。在 \(p\)-\(V\) 图上，哪个过程的终点压强更大？哪个消耗的功更多？（思路提示：绝热过程不散热，压缩功全部变成内能，温度升高，压强因此比等温更高；绝热线更陡，终点在等温线上方，曲线下面积更大，耗功更多。这也解释了打气筒为什么越打越费劲——不全是轮胎压强高了，还有温度在帮倒忙。）
  \item 用手堵住打气筒的出口，快速把活塞压到底再松开，活塞会被气体顶回一段距离，但回不到原来的高度。这笔能量账怎么平？（思路提示：压缩时你对气体做功，气体内能升高；回弹时气体反过来对你做功，内能回落。但压缩与回弹都很快，气体内部产生扰动和涡流，加上筒壁散热，一部分功变成了散不回来的内能。回来的功少于进去的功，差额留在气体和筒壁里——机械能账上看似“消失”，总账依然平衡。再注意：这个过程是单向的，把录像倒放会一眼看出假来。为什么账能平、录像却不能倒放？这是本章留给第 29 章的钩子。）
\end{puzzles}
